\subsection{The Cross Product}

\content{
\begin{definition} (The Cross Product) If $\vec{a}=\la a_1,a_2,a_3\ra$ and 
$\vec{b} = \la b_1,b_2,b_3\ra$, then the cross product, denoted $\vec{a}\times
\vec{b}$ is given by 
$$ \vec{a}\times \vec{b} = \la a_2b_3-a_3b_2,a_3b_1 - a_1b_3,a_1b_2-a_2b_1\ra.$$
\end{definition}
}{% presentation
\vspace{2in}
}{% nopresentation
\vspace{1in}
}{% comments
Show how this can be remembered by taking the determinant of a 3x3 matrix. 
}

\content{
\begin{example} Find the cross product between the vectors 
$\vec{a}=\la 1,3,4\ra$ and $\vec{b}=\la 2,7,-5\ra$.
\end{example}
}{% presentation
\vspace{2in}
}{% nopresentation
\vspace{1in}
}{% comments
$\vec{a}\times\vec{b}=\la -43, 13, 1\ra$.
}

\content{
\begin{theorem} The vector $\vec{a}\times\vec{b}$ is orthogonal to both
$\vec{a}$ and $\vec{b}$.
\end{theorem}
}{% presentation
}{% nopresentation
}{% comments
}

\content{
\begin{theorem} If $\theta$ is the angle between vectors $\vec{a}$ and
$\vec{b}$, then 
$$ |\vec{a}\times\vec{b}| = |\vec{a}||\vec{b}|\sin\theta. $$
\end{theorem}
}{% presentation
\vspace{2in}
}{% nopresentation
\vspace{1in}
}{% comments
Use the space below to show that this is precesily the area of the parallelogram
spanned by $\vec{a}$ and $\vec{b}$, so that the magnitude of the cross product
gives us the area of that parallelogram. 
}

\content{
\begin{theorem} Two vectors $\vec{a}$ and $\vec{b}$ are parallel if and only if 
$\vec{a}\times\vec{b}=0$. 
\end{theorem}
}{% presentation
}{% nopresentation
}{% comments
}

\content{
\begin{example} Find a vector which is perpendicular to the plane passing
through the points $P(1,4,6)$, $Q(-2,5,-1)$, and $R(1,-1,1)$.
\end{example}
}{% presentation
\vspace{3in}
}{% nopresentation
\vspace{2in}
\newpage
}{% comments
One answer is $\la -40, -15, 15\ra$.
}

\content{
\begin{theorem} (Properties of the Cross Product) If $\vec{a}$, $\vec{b}$, and 
$\vec{c}$ are vectors and $c$ is a scalar, then 
\begin{enumerate}
\item $\vec{a}\times\vec{b} = -\vec{b}\times\vec{a}$
\item $\vec{a}\times (\vec{b}+\vec{c}) =
\vec{a}\times\vec{b}+\vec{a}\times\vec{c}$
\item $(c\vec{a})\times\vec{b}=c(\vec{a}\times\vec{b}) = 
\vec{a}\times(c\vec{b})$
\item $(\vec{a}+\vec{b})\times \vec{c} = \vec{a}\times \vec{c} +
\vec{b}\times\vec{c}$
\item $\vec{a}\cdot (\vec{b}\times\vec{c}) = 
(\vec{a}\times\vec{b})\cdot\vec{c}$
\item $\vec{a}\times(\vec{b}\times\vec{c})=
(\vec{a}\cdot\vec{c})\vec{b}-(\vec{a}\cdot \vec{b})\vec{c}$
\end{enumerate}
\end{theorem}
}{% presentation
\vspace{4in}
}{% nopresentation
\vspace{2in}
}{% comments
Maybe give a proof of one of these. 
}

\content{
The product $\vec{a}\cdot (\vec{b}\times\vec{c})$ that occurs in the above
theorem is known as the scalar triple product. And it's value is the volume of
the parallelpiped determined by the vectors $\vec{a}$, $\vec{b}$, and $\vec{c}$. 
}{% presentation
\vspace{2in}
}{% nopresentation
\vspace{1in}
}{% comments
Don't give a proof, just draw out a parallelpiped, and move onto the next
example. 
}

\content{
\begin{example} Find the volume of the parallelpiped determined by the three
vectors $\vec{a}=\la 1,4,-7\ra$, $\vec{b}=\la 2,-1,4\ra$, and 
$\vec{c} =\la 0,-9,18\ra$. 
\end{example}
}{% presentation
\vspace{2in}
}{% nopresentation
\vspace{1.5in}
}{% comments
The volume is 0, which means that these vectors actually lie in the same plane,
that is, they are coplanar. 

An Application to computing torque is outlined in the textbook for those who are
interested in reading that. 
}

\content{
If we have a force $\vec{F}$ acting on a rigid body at a point
given a position vector $\vec{r}$, the torque is defined $\vec{\tau}$ is defined
to be the cross product of $\vec{r}$ and $\vec{F}$:
$$ \vec{\tau} = \vec{r}\times \vec{F}. $$
The torque is a measure of the tendancy of the body to rotate, and the direction
of the torque vector is the axis of rotation (rotating clockwise about this
axis)

\begin{example}
A bolt is tightened by applying a 40 N force to a 0.25 m wrench. Find the
magnitude of the torque about the center of the bolt.
\end{example}
}{% presentation
\vspace{2in}
}{% nopresentation
\vspace{1in}
}{% comments
}
